\section{Especificación de Historias de Usuario}

Esta sección detalla las Historias de Usuario que sirvieron como base para la elicitación de los requisitos del sistema. Cada historia representa una unidad de valor funcional desde la perspectiva de un actor específico.

\begin{quote}
\textbf{Nota sobre Nomenclatura:} Se ha implementado un sistema de identificación global único para garantizar la trazabilidad. Cada historia posee un código compuesto por el Prefijo del Módulo (MTTO, INV, VIS, ADM) seguido de un consecutivo numérico único en todo el proyecto.
\end{quote}

\begin{quote}
\textbf{Nota sobre la Estimación de Esfuerzo:}

Para la valoración de la complejidad de cada historia, se ha utilizado la escala de \textbf{Fibonacci} estándar en metodologías ágiles (Cohn, 2004). Los puntos reflejan una combinación de esfuerzo, complejidad técnica y riesgo asociado:

\begin{itemize}
    \item \textbf{1 - 2 Puntos:} Tareas triviales o de muy baja complejidad (ajustes menores).
    \item \textbf{3 Puntos:} Funcionalidades estándar con baja incertidumbre (ej. CRUDs básicos, formularios).
    \item \textbf{5 Puntos:} Funcionalidades con lógica de negocio compleja o integraciones con sistemas externos.
    \item \textbf{8 Puntos:} Alta complejidad técnica, requerimientos no funcionales críticos (Offline, Seguridad) o alta incertidumbre (Investigación requerida).
\end{itemize}
\end{quote}

A continuación, se presentan las fichas técnicas agrupadas por módulo funcional.

\subsection{Módulo de Gestión de Operaciones (MTTO)}

\noindent
\textbf{Ficha Técnica: MTTO-001 - Programación de Mantenimiento Preventivo}
\footnotesize
\begin{longtable}{|p{0.3\textwidth}|p{0.65\textwidth}|}
\hline
\textbf{Usuario (Rol)} & Planificador \\ \hline
\textbf{Puntos estimados} & 5 \\ \hline
\textbf{Descripción} & \textbf{Como} Planificador, \textbf{Quiero} crear y programar Órdenes de Trabajo (WO) preventivas de forma recurrente (ej. mensual, trimestral), \textbf{Para} asegurar que el mantenimiento rutinario se ejecute a tiempo sin intervención manual, reduciendo el riesgo de fallos por olvido y optimizando la confiabilidad de los activos. \\ \hline
\textbf{Observaciones} & Requiere un motor de tareas programadas (Cron jobs) en el backend. \\ \hline
\end{longtable}
\normalsize

\noindent
\textbf{Criterios de Aceptación:}
\footnotesize
\begin{longtable}{|p{0.25\textwidth}|p{0.2\textwidth}|p{0.2\textwidth}|p{0.25\textwidth}|}
\hline
\textbf{Escenario} & \textbf{Dado} & \textbf{Cuando} & \textbf{Entonces} \\ \hline
\endhead
Creación exitosa de una Orden de Trabajo (WO) Preventiva & Dado que el `Planificador' está autenticado en el sistema y se encuentra en la pantalla de ``Planificación de Mantenimiento Preventivo''. & Cuando selecciona un activo, define una lista de tareas, establece una frecuencia (ej. ``mensual''), y guarda el plan. & Entonces el sistema debe confirmar la creación del plan recurrente y la próxima WO debe aparecer programada en el calendario de mantenimiento. \\ \hline
Intento de creación con datos incompletos & Dado que el `Planificador' está autenticado en el sistema y se encuentra en la pantalla de ``Planificación de Mantenimiento Preventivo''. & Cuando intenta guardar un plan sin haber seleccionado un activo o sin definir una frecuencia. & Entonces el sistema debe mostrar un mensaje de error claro indicando los campos obligatorios que faltan y no debe crear la orden. \\ \hline
\end{longtable}
\normalsize

\noindent
\textbf{Ficha Técnica: MTTO-002 - Cierre Digital de OT en Móvil}
\footnotesize
\begin{longtable}{|p{0.3\textwidth}|p{0.65\textwidth}|}
\hline
\textbf{Usuario (Rol)} & Técnico \\ \hline
\textbf{Puntos estimados} & 8 \\ \hline
\textbf{Descripción} & \textbf{Como} Técnico, \textbf{Quiero} registrar el cierre de una orden de trabajo desde mi tablet, llenando un formulario digital y adjuntando fotos/videos, \textbf{Para} documentar la reparación en el momento, sin tener que volver a la oficina a transcribir notas en papel.. \\ \hline
\textbf{Observaciones} & Complejidad alta debido al requisito de funcionamiento Offline y sincronización. \\ \hline
\end{longtable}
\normalsize

\noindent
\textbf{Criterios de Aceptación:}
\footnotesize
\begin{longtable}{|p{0.25\textwidth}|p{0.2\textwidth}|p{0.2\textwidth}|p{0.25\textwidth}|}
\hline
\textbf{Escenario} & \textbf{Dado} & \textbf{Cuando} & \textbf{Entonces} \\ \hline
\endhead
Cierre Exitoso de una WO en Campo & Dado que el `Técnico' ha finalizado una reparación y está en la pantalla de la WO en su tablet. & Cuando completa los campos obligatorios del formulario de cierre (ej. horas, causa de falla) y presiona ``Finalizar''. & Entonces el estado de la WO debe cambiar a ``Pendiente de Revisión'', la información debe guardarse en el historial del activo y el `Supervisor' debe recibir una notificación. \\ \hline
Adjuntar Evidencia Multimedia & Dado que el `Técnico' está llenando el formulario de cierre en su tablet. & Cuando utiliza la función ``Adjuntar Foto/Video'' y captura una imagen de la reparación finalizada. & Entonces la imagen debe quedar asociada a la orden de trabajo y ser visible para el Supervisor. \\ \hline
Intento de Cierre con Datos Incompletos & Dado que el `Técnico' está en el formulario de cierre. & Cuando intenta finalizar la WO sin haber llenado un campo obligatorio (ej. ``causa de la falla''). & Entonces el sistema debe mostrar un mensaje de error claro, resaltar el campo faltante y \textbf{no permitir} el cierre de la orden. \\ \hline
Operación sin Conexión (Offline) & Dado que el `Técnico' está en una zona de la planta sin WiFi y ha abierto una WO. & Cuando completa el formulario de cierre y presiona ``Finalizar''. & Entonces la aplicación debe guardar los datos localmente en la tablet y mostrar un indicador de ``Pendiente de Sincronización''. \\ \hline
\end{longtable}
\normalsize

\noindent
\textbf{Ficha Técnica: MTTO-003 - Carga de Informe por Contratista}
\footnotesize
\begin{longtable}{|p{0.3\textwidth}|p{0.65\textwidth}|}
\hline
\textbf{Usuario (Rol)} & Contratista \\ \hline
\textbf{Puntos estimados} & 3 \\ \hline
\textbf{Descripción} & \textbf{Como} Contratista, \textbf{Quiero} subir mi informe de trabajo en formato PDF a la orden de trabajo asignada, \textbf{Para} cumplir con el requisito de entrega del cliente de forma rápida y digital.. \\ \hline
\textbf{Observaciones} & Requiere validación de tipos de archivo (MIME types) y escaneo de seguridad. \\ \hline
\end{longtable}
\normalsize

\noindent
\textbf{Criterios de Aceptación:}
\footnotesize
\begin{longtable}{|p{0.25\textwidth}|p{0.2\textwidth}|p{0.2\textwidth}|p{0.25\textwidth}|}
\hline
\textbf{Escenario} & \textbf{Dado} & \textbf{Cuando} & \textbf{Entonces} \\ \hline
\endhead
Carga exitosa de informe PDF & Dado que el `Contratista' está autenticado y en la pantalla de la WO. & Cuando selecciona un archivo PDF válido y hace clic en ``Subir''. & Entonces el sistema debe confirmar que el archivo se ha adjuntado a la WO, mostrar el nombre del archivo y enviar una notificación al `Supervisor'. \\ \hline
Intento de carga con formato incorrecto & Dado que el `Contratista' está autenticado y en la pantalla de la WO. & Cuando intenta subir un archivo que \textbf{no es} PDF o de imagen (ej. un `.exe'). & Entonces el sistema debe mostrar un mensaje de error claro (``Formato de archivo no válido'') y \textbf{no debe} adjuntar el archivo. \\ \hline
\end{longtable}
\normalsize

\noindent
\textbf{Ficha Técnica: MTTO-004 - Extracción Automática de Datos de Informes}
\footnotesize
\begin{longtable}{|p{0.3\textwidth}|p{0.65\textwidth}|}
\hline
\textbf{Usuario (Rol)} & Supervisor \\ \hline
\textbf{Puntos estimados} & 8 \\ \hline
\textbf{Descripción} & \textbf{Como} Supervisor, \textbf{Quiero} que el sistema extraiga automáticamente la información clave (horas, piezas, descripción) del PDF que subió el contratista y me la presente en un formato para aprobar con un solo clic, \textbf{Para} ahorrar el tiempo que hoy pierdo leyendo y transcribiendo manualmente cada informe para cerrar una orden.. \\ \hline
\textbf{Observaciones} & Dependencia crítica de la integración con el servicio de IA (LLM). Riesgo técnico alto. \\ \hline
\end{longtable}
\normalsize

\noindent
\textbf{Criterios de Aceptación:}
\footnotesize
\begin{longtable}{|p{0.25\textwidth}|p{0.2\textwidth}|p{0.2\textwidth}|p{0.25\textwidth}|}
\hline
\textbf{Escenario} & \textbf{Dado} & \textbf{Cuando} & \textbf{Entonces} \\ \hline
\endhead
Extracción de datos exitosa & Dado que un `Contratista' ha subido un PDF a una WO y está pendiente de aprobación. & Cuando el `Supervisor' abre la WO. & Entonces el sistema debe mostrarle un resumen con los datos clave (horas, piezas) ya extraídos del PDF por la IA. \\ \hline
Extracción con datos faltantes & Dado que un `Contratista' ha subido un PDF que no contiene la ``lista de materiales''. & Cuando el `Supervisor' abre la WO. & Entonces el sistema debe mostrar un \textbf{aviso de validación} que diga ``Información Faltante: No se pudo extraer la lista de materiales. Por favor, revise el documento manualmente.'' \\ \hline
\end{longtable}
\normalsize

\noindent
\textbf{Ficha Técnica: MTTO-020 - Consulta de Manuales en Móvil}
\footnotesize
\begin{longtable}{|p{0.3\textwidth}|p{0.65\textwidth}|}
\hline
\textbf{Usuario (Rol)} & Técnico \\ \hline
\textbf{Puntos estimados} & 3 \\ \hline
\textbf{Descripción} & \textbf{Como} Técnico, \textbf{Quiero} buscar y descargar los manuales técnicos y planos del activo directamente desde la aplicación móvil, \textbf{Para} consultar la información necesaria en el sitio de la reparación sin tener que desplazarme a la oficina o buscar carpetas físicas.. \\ \hline
\textbf{Observaciones} & Requiere visor de PDF integrado en la aplicación móvil. \\ \hline
\end{longtable}
\normalsize

\noindent
\textbf{Criterios de Aceptación:}
\footnotesize
\begin{longtable}{|p{0.25\textwidth}|p{0.2\textwidth}|p{0.2\textwidth}|p{0.25\textwidth}|}
\hline
\textbf{Escenario} & \textbf{Dado} & \textbf{Cuando} & \textbf{Entonces} \\ \hline
\endhead
\textbf{Visualización de Documento Técnico} & Dado que el Técnico ha escaneado el QR de una bomba y está en su ficha móvil. & Cuando selecciona la pestaña ``Documentación'' y toca el archivo ``Manual de Servicio.pdf''. & Entonces el documento debe abrirse en el visor integrado de la tablet sin necesidad de salir de la aplicación. \\ \hline
\textbf{Descarga para uso Offline} & Dado que el Técnico tiene conexión a internet y visualiza la lista de planos. & Cuando presiona el botón ``Descargar'' o ``Hacer disponible sin conexión'' sobre un plano eléctrico. & Entonces el archivo debe guardarse localmente y un icono debe indicar que está disponible para futuras consultas sin internet. \\ \hline
\end{longtable}
\normalsize

\subsection{Módulo de Gestión de Recursos (INV)}

\noindent
\textbf{Ficha Técnica: INV-005 - Creación de Ficha Técnica de Activo}
\footnotesize
\begin{longtable}{|p{0.3\textwidth}|p{0.65\textwidth}|}
\hline
\textbf{Usuario (Rol)} & Ingeniero de Proyectos \\ \hline
\textbf{Puntos estimados} & 3 \\ \hline
\textbf{Descripción} & \textbf{Como} Ingeniero de Proyectos, \textbf{Quiero} crear la ficha técnica digital de un nuevo activo (Nombre, Serial, Fecha de Compra, Garantía), \textbf{Para} tener un registro único oficial en el sistema antes de asociarle modelos 3D o planes de mantenimiento.. \\ \hline
\textbf{Observaciones} & N/A \\ \hline
\end{longtable}
\normalsize

\noindent
\textbf{Criterios de Aceptación:}
\footnotesize
\begin{longtable}{|p{0.25\textwidth}|p{0.2\textwidth}|p{0.2\textwidth}|p{0.25\textwidth}|}
\hline
\textbf{Escenario} & \textbf{Dado} & \textbf{Cuando} & \textbf{Entonces} \\ \hline
\endhead
Alta de Activo Nuevo & Dado que el Ingeniero está en el módulo de ``Inventario de Activos''. & Cuando hace clic en ``Nuevo'', llena los datos obligatorios (Serial, Modelo) y guarda. & Entonces el sistema debe crear el activo con un estado ``Activo'' y generar un código QR único para su identificación. \\ \hline
Validación de Serial Único & Dado que ya existe una bomba con el serial ``P-101'' en la base de datos. & Cuando el Ingeniero intenta crear otro activo con el mismo serial ``P-101''. & Entonces el sistema debe impedir el guardado y mostrar el error: ``El serial ya está registrado en el sistema''. \\ \hline
\end{longtable}
\normalsize

\noindent
\textbf{Ficha Técnica: INV-006 - Ajuste Manual de Stock de Repuestos}
\footnotesize
\begin{longtable}{|p{0.3\textwidth}|p{0.65\textwidth}|}
\hline
\textbf{Usuario (Rol)} & Jefe de Almacén \\ \hline
\textbf{Puntos estimados} & 3 \\ \hline
\textbf{Descripción} & \textbf{Como} Jefe de Almacén, \textbf{Quiero} ajustar manualmente el stock de un repuesto (entrada por compra o ajuste de inventario), \textbf{Para} mantener la consistencia entre el inventario físico real y lo que muestra el sistema.. \\ \hline
\textbf{Observaciones} & Debe generar un registro de auditoría (Log) inmutable automáticamente. \\ \hline
\end{longtable}
\normalsize

\noindent
\textbf{Criterios de Aceptación:}
\footnotesize
\begin{longtable}{|p{0.25\textwidth}|p{0.2\textwidth}|p{0.2\textwidth}|p{0.25\textwidth}|}
\hline
\textbf{Escenario} & \textbf{Dado} & \textbf{Cuando} & \textbf{Entonces} \\ \hline
\endhead
Entrada de Mercancía & Dado que el Jefe de Almacén tiene 5 unidades del repuesto ``Rodamiento SKF'' en sistema. & Cuando registra una entrada de compra de 10 unidades nuevas. & Entonces el stock disponible debe actualizarse automáticamente a 15 unidades y quedar registro del movimiento. \\ \hline
\end{longtable}
\normalsize

\noindent
\textbf{Ficha Técnica: INV-007 - Validación y Alta de Activos Nuevos}
\footnotesize
\begin{longtable}{|p{0.3\textwidth}|p{0.65\textwidth}|}
\hline
\textbf{Usuario (Rol)} & Ingeniero de Proyectos \\ \hline
\textbf{Puntos estimados} & 3 \\ \hline
\textbf{Descripción} & \textbf{Como} Ingeniero de Proyectos, \textbf{Quiero} revisar y completar la ficha técnica de los activos nuevos que están en estado ``Borrador'' o ``En Cuarentena'', \textbf{Para} asegurar que todos los datos técnicos (manuales, modelo 3D, parámetros) sean correctos antes de liberar el activo para su operación y mantenimiento.. \\ \hline
\textbf{Observaciones} & Flujo de cambio de estado. Requiere validación de campos obligatorios antes de activar. \\ \hline
\end{longtable}
\normalsize

\noindent
\textbf{Criterios de Aceptación:}
\footnotesize
\begin{longtable}{|p{0.25\textwidth}|p{0.2\textwidth}|p{0.2\textwidth}|p{0.25\textwidth}|}
\hline
\textbf{Escenario} & \textbf{Dado} & \textbf{Cuando} & \textbf{Entonces} \\ \hline
\endhead
Validación y Activación de Activo & Dado que existe un registro ``Bomba P-New'' creado por Almacén en estado ``Borrador''. & Cuando el Ingeniero completa los campos técnicos obligatorios y cambia el estado a ``Operativo''. & Entonces el activo debe ser visible para los Técnicos en el mapa y disponible para asignarle planes de mantenimiento. \\ \hline
Rechazo de Activo Incompleto & Dado que el Ingeniero revisa un activo en ``Borrador'' que no corresponde a la compra. & Cuando selecciona la opción ``Rechazar/Devolver''. & Entonces el registro debe marcarse como ``Devuelto a Almacén'' y se debe notificar al Jefe de Almacén para que revise la entrada física. \\ \hline
\end{longtable}
\normalsize

\noindent
\textbf{Ficha Técnica: INV-021 - Alertas de Stock Crítico}
\footnotesize
\begin{longtable}{|p{0.3\textwidth}|p{0.65\textwidth}|}
\hline
\textbf{Usuario (Rol)} & Jefe de Almacén \\ \hline
\textbf{Puntos estimados} & 5 \\ \hline
\textbf{Descripción} & \textbf{Como} Jefe de Almacén, \textbf{Quiero} recibir una notificación automática cuando el stock de un ``Repuesto Crítico'' baje del mínimo establecido, \textbf{Para} reponer el inventario a tiempo y garantizar la disponibilidad de piezas esenciales para emergencias. \\ \hline
\textbf{Observaciones} & Proceso en segundo plano (Background Worker) para monitorear niveles de stock. \\ \hline
\end{longtable}
\normalsize

\noindent
\textbf{Criterios de Aceptación:}
\footnotesize
\begin{longtable}{|p{0.25\textwidth}|p{0.2\textwidth}|p{0.2\textwidth}|p{0.25\textwidth}|}
\hline
\textbf{Escenario} & \textbf{Dado} & \textbf{Cuando} & \textbf{Entonces} \\ \hline
\endhead
\textbf{Disparo de Alerta por Consumo} & Dado que el repuesto ``Sello Mecánico'' tiene un stock mínimo definido de 5 unidades y un stock actual de 5. & Cuando el sistema procesa una salida de almacén de 1 unidad (dejando el stock en 4). & Entonces el sistema debe generar inmediatamente una notificación interna y enviar un correo al Jefe de Almacén con el asunto ``Alerta: Stock Crítico - Sello Mecánico''. \\ \hline
\textbf{Visualización de Ítems Críticos} & Dado que existen 3 repuestos por debajo del mínimo. & Cuando el Jefe de Almacén entra al Dashboard de Inventario. & Entonces debe ver un widget o tarjeta roja destacada que liste esos 3 ítems con la etiqueta ``Reabastecimiento Urgente''. \\ \hline
\end{longtable}
\normalsize

\subsection{Módulo de Visualización y Gemelo Digital (VIS)}

\noindent
\textbf{Ficha Técnica: VIS-008 - Visualización de Capa de Permisos de Trabajo}
\footnotesize
\begin{longtable}{|p{0.3\textwidth}|p{0.65\textwidth}|}
\hline
\textbf{Usuario (Rol)} & Inspector HSEQ \\ \hline
\textbf{Puntos estimados} & 5 \\ \hline
\textbf{Descripción} & \textbf{Como} Inspector HSEQ, \textbf{Quiero} habilitar una capa visual que muestre los permisos de trabajo activos (alturas, caliente, izaje) sobre el mapa de la planta, \textbf{Para} identificar rápidamente cruces de actividades peligrosas y evitar accidentes por interferencia entre contratistas.. \\ \hline
\textbf{Observaciones} & Requiere manejo de coordenadas geoespaciales sobre el plano de planta. \\ \hline
\end{longtable}
\normalsize

\noindent
\textbf{Criterios de Aceptación:}
\footnotesize
\begin{longtable}{|p{0.25\textwidth}|p{0.2\textwidth}|p{0.2\textwidth}|p{0.25\textwidth}|}
\hline
\textbf{Escenario} & \textbf{Dado} & \textbf{Cuando} & \textbf{Entonces} \\ \hline
\endhead
Activación de la Capa de Seguridad & Dado que el `Inspector' HSEQ está visualizando el mapa general de la planta en 2D o 3D. & Cuando hace clic en el botón ``Capas'' y selecciona la opción ``Permisos de Trabajo''. & Entonces el mapa debe superponer iconos específicos en las coordenadas donde hay trabajos activos, diferenciándolos por color según el tipo de riesgo (ej. Rojo para Caliente, Azul para Alturas). \\ \hline
Consulta de Detalle de Permiso & Dado que la capa de ``Permisos de Trabajo'' está activa y se muestran iconos en el mapa. & Cuando el Inspector hace clic sobre un icono de ``Trabajo en Alturas''. & Entonces se debe desplegar una tarjeta emergente mostrando: el nombre del contratista, la hora de inicio/fin del permiso y el responsable a cargo. \\ \hline
\end{longtable}
\normalsize

\noindent
\textbf{Ficha Técnica: VIS-009 - Navegación con Zoom en Mapa de Planta}
\footnotesize
\begin{longtable}{|p{0.3\textwidth}|p{0.65\textwidth}|}
\hline
\textbf{Usuario (Rol)} & Supervisor \\ \hline
\textbf{Puntos estimados} & 3 \\ \hline
\textbf{Descripción} & \textbf{Como} Supervisor, \textbf{Quiero} hacer zoom en una zona específica de la planta usando el scroll del mouse, \textbf{Para} inspeccionar con detalle la ubicación de los activos en áreas densas. \\ \hline
\textbf{Observaciones} & Se debe asegurar la compatibilidad con gestos táctiles (Touch events). \\ \hline
\end{longtable}
\normalsize

\noindent
\textbf{Criterios de Aceptación:}
\footnotesize
\begin{longtable}{|p{0.25\textwidth}|p{0.2\textwidth}|p{0.2\textwidth}|p{0.25\textwidth}|}
\hline
\textbf{Escenario} & \textbf{Dado} & \textbf{Cuando} & \textbf{Entonces} \\ \hline
\endhead
Zoom In exitoso & Dado que el `Supervisor' está visualizando el mapa general de la planta en 2D. & Cuando gira la rueda del mouse hacia arriba sobre la zona de ``Calderas'' & Entonces el mapa debe acercarse centrado en el cursor y mostrar los iconos de los activos más pequeños \\ \hline
Límite de Zoom & Dado que el `Supervisor' ha alcanzado el nivel máximo de acercamiento. & Cuando intenta hacer más zoom. & Entonces la vista no debe cambiar y no debe pixelarse la imagen de fondo. \\ \hline
\end{longtable}
\normalsize

\noindent
\textbf{Ficha Técnica: VIS-010 - Visualización de Despiece (Vista Explosionada)}
\footnotesize
\begin{longtable}{|p{0.3\textwidth}|p{0.65\textwidth}|}
\hline
\textbf{Usuario (Rol)} & Técnico \\ \hline
\textbf{Puntos estimados} & 8 \\ \hline
\textbf{Descripción} & \textbf{Como} Técnico, \textbf{Quiero} visualizar la vista explosionada (despiece) del activo seleccionado en 3D, \textbf{Para} identificar la ubicación exacta de un componente interno antes de iniciar el desarme físico.. \\ \hline
\textbf{Observaciones} & Complejidad alta en la manipulación del grafo de escena 3D para la animación de despiece. \\ \hline
\end{longtable}
\normalsize

\noindent
\textbf{Criterios de Aceptación:}
\footnotesize
\begin{longtable}{|p{0.25\textwidth}|p{0.2\textwidth}|p{0.2\textwidth}|p{0.25\textwidth}|}
\hline
\textbf{Escenario} & \textbf{Dado} & \textbf{Cuando} & \textbf{Entonces} \\ \hline
\endhead
Activación de Despiece & Dado que el `Técnico' ha seleccionado una bomba y ve el menú de opciones del activo. & Cuando selecciona la opción ``Ver Despiece''. & Entonces las partes del modelo 3D deben separarse visualmente (animación de explosión) permitiendo ver los componentes internos individualmente. \\ \hline
Selección de Componente Interno & Dado que el activo se encuentra en modo ``Vista Explosionada''. & Cuando el Técnico selecciona un componente específico (ej. rodamiento). & Entonces el sistema debe resaltar ese componente y mostrar una etiqueta con su nombre y código de parte. \\ \hline
\end{longtable}
\normalsize

\noindent
\textbf{Ficha Técnica: VIS-011 - Visualización de Rutas de Aislamiento (LOTO)}
\footnotesize
\begin{longtable}{|p{0.3\textwidth}|p{0.65\textwidth}|}
\hline
\textbf{Usuario (Rol)} & Técnico \\ \hline
\textbf{Puntos estimados} & 8 \\ \hline
\textbf{Descripción} & \textbf{Como} Técnico, \textbf{Quiero} visualizar la ruta de aislamiento de energía (LOTO) sobre el modelo 3D, \textbf{Para} identificar rápidamente los puntos de bloqueo eléctrico/mecánico y evitar accidentes.. \\ \hline
\textbf{Observaciones} & Requiere lógica de trazado de rutas entre nodos 3D. \\ \hline
\end{longtable}
\normalsize

\noindent
\textbf{Criterios de Aceptación:}
\footnotesize
\begin{longtable}{|p{0.25\textwidth}|p{0.2\textwidth}|p{0.2\textwidth}|p{0.25\textwidth}|}
\hline
\textbf{Escenario} & \textbf{Dado} & \textbf{Cuando} & \textbf{Entonces} \\ \hline
\endhead
Visualización de Puntos LOTO & Dado que el `Técnico' tiene seleccionado un activo en la tablet. & Cuando activa la opción ``Ver Puntos de Aislamiento (LOTO)''. & Entonces la cámara debe alejarse para mostrar el contexto y se deben resaltar con una línea brillante las conexiones hacia los tableros eléctricos o válvulas de corte asociados. \\ \hline
Validación de Conexión & Dado que se muestran las líneas de conexión LOTO. & Cuando el Técnico selecciona el punto de corte sugerido (ej. Breaker). & Entonces el sistema debe mostrar la etiqueta del tablero (Tag) que coincide con la realidad física. \\ \hline
\end{longtable}
\normalsize

\noindent
\textbf{Ficha Técnica: VIS-012 - Configuración de Jerarquía de Componentes 3D}
\footnotesize
\begin{longtable}{|p{0.3\textwidth}|p{0.65\textwidth}|}
\hline
\textbf{Usuario (Rol)} & Ingeniero de Proyectos \\ \hline
\textbf{Puntos estimados} & 5 \\ \hline
\textbf{Descripción} & \textbf{Como} Ingeniero de Proyectos, \textbf{Quiero} definir la jerarquía de componentes de un activo (ej. Bomba $\to$ Motor $\to$ Rodamiento) y asociar un modelo 3D a cada nivel, \textbf{Para} construir la vista explosionada del Gemelo Digital y permitir que los técnicos accedan a la información del componente específico en campo.. \\ \hline
\textbf{Observaciones} & Estructura de datos recursiva (Árbol). Requiere validación de ciclos. \\ \hline
\end{longtable}
\normalsize

\noindent
\textbf{Criterios de Aceptación:}
\footnotesize
\begin{longtable}{|p{0.25\textwidth}|p{0.2\textwidth}|p{0.2\textwidth}|p{0.25\textwidth}|}
\hline
\textbf{Escenario} & \textbf{Dado} & \textbf{Cuando} & \textbf{Entonces} \\ \hline
\endhead
Crear un sub-componente & Dado que el Ingeniero está en la página de un activo ``Bomba P-101'' que no tiene componentes definidos. & Cuando hace clic en ``Añadir Componente'', le da el nombre ``Motor Eléctrico'' y el código ``ME-01''. & Entonces ``Motor Eléctrico'' debe aparecer como un sub-ítem anidado debajo de la ``Bomba P-101'' en el árbol de activos. \\ \hline
Asociar un modelo 3D a un componente & Dado que el Ingeniero está editando las propiedades del componente ``Motor Eléctrico''. & Cuando utiliza la opción ``Subir Modelo 3D'' y adjunta un archivo `motor.step'. & Entonces el sistema debe asociar ese archivo 3D al componente y mostrar una vista previa en miniatura. \\ \hline
Navegación jerárquica & Dado que el Técnico está viendo el activo ``Bomba P-101'' en la app móvil. & Cuando hace clic en la vista explosionada y selecciona el ``Motor Eléctrico''. & Entonces la app debe mostrar la información y el modelo 3D específico del ``Motor Eléctrico'', no el de la bomba completa. \\ \hline
Intento de crear un componente duplicado & Dado que el activo ``Bomba P-101'' ya tiene un componente con el código ``ME-01''. & Cuando el Ingeniero intenta crear otro componente con el mismo código ``ME-01''. & Entonces el sistema debe mostrar un mensaje de error ``El código de componente ya existe en este activo'' y no debe crear el duplicado. \\ \hline
Mover componente en el árbol (Re-parenting) & Dado que el componente ``Rodamiento R-01'' fue creado por error dentro del ``Motor A''. & Cuando el Ingeniero arrastra el componente hacia el ``Motor B'' en el árbol de jerarquía. & Entonces el sistema debe actualizar la relación padre-hijo, moviendo el rodamiento y todo su historial al nuevo padre. \\ \hline
\end{longtable}
\normalsize

\subsection{Módulo de Administración e Inteligencia (ADM)}

\noindent
\textbf{Ficha Técnica: ADM-013 - Gestión de Roles y Permisos}
\footnotesize
\begin{longtable}{|p{0.3\textwidth}|p{0.65\textwidth}|}
\hline
\textbf{Usuario (Rol)} & Administrador \\ \hline
\textbf{Puntos estimados} & 5 \\ \hline
\textbf{Descripción} & \textbf{Como} Administrador, \textbf{Quiero} definir y gestionar los roles de usuario (ej. Técnico, Supervisor), \textbf{Para} controlar qué conjunto de acciones puede realizar cada tipo de usuario en el sistema.. \\ \hline
\textbf{Observaciones} & Impacto crítico en la seguridad. Requiere ACL (Access Control List). \\ \hline
\end{longtable}
\normalsize

\noindent
\textbf{Criterios de Aceptación:}
\footnotesize
\begin{longtable}{|p{0.25\textwidth}|p{0.2\textwidth}|p{0.2\textwidth}|p{0.25\textwidth}|}
\hline
\textbf{Escenario} & \textbf{Dado} & \textbf{Cuando} & \textbf{Entonces} \\ \hline
\endhead
Asignar permisos a un Rol & Dado que el `Administrador' está en la pantalla de ``Gestión de Roles'' y ha seleccionado el rol ``Técnico''. & Cuando marca el checkbox del permiso ``Crear Orden de Trabajo'' y guarda los cambios. & Entonces un usuario con el rol ``Técnico'' ahora debe ver el botón para crear órdenes de trabajo. \\ \hline
\end{longtable}
\normalsize

\noindent
\textbf{Ficha Técnica: ADM-014 - Gestión de Usuarios y Acceso}
\footnotesize
\begin{longtable}{|p{0.3\textwidth}|p{0.65\textwidth}|}
\hline
\textbf{Usuario (Rol)} & Administrador \\ \hline
\textbf{Puntos estimados} & 3 \\ \hline
\textbf{Descripción} & \textbf{Como} Administrador, \textbf{Quiero} crear una cuenta para un nuevo empleado y asignarle un rol existente, \textbf{Para} darle acceso seguro al sistema según sus responsabilidades.. \\ \hline
\textbf{Observaciones} & Integración con servicio de correo para notificaciones de bienvenida. \\ \hline
\end{longtable}
\normalsize

\noindent
\textbf{Criterios de Aceptación:}
\footnotesize
\begin{longtable}{|p{0.25\textwidth}|p{0.2\textwidth}|p{0.2\textwidth}|p{0.25\textwidth}|}
\hline
\textbf{Escenario} & \textbf{Dado} & \textbf{Cuando} & \textbf{Entonces} \\ \hline
\endhead
Creación exitosa de Usuario & Dado que el `Administrador' está en la pantalla de ``Gestión de Usuarios''. & Cuando ingresa el nombre, correo y selecciona el rol ``Técnico'' para un nuevo usuario y guarda. & Entonces el sistema debe crear la cuenta y enviar un correo de bienvenida a esa dirección. \\ \hline
Intento de crear usuario duplicado & Dado que ya existe un usuario con el correo ``juliana@dtmtto.com''. & Cuando el `Administrador' intenta crear un nuevo usuario con ese mismo correo. & Entonces el sistema debe mostrar un mensaje de error ``El correo ya está en uso'' y no debe crear la cuenta. \\ \hline
\end{longtable}
\normalsize

\noindent
\textbf{Ficha Técnica: ADM-015 - Configuración de Integración ERP}
\footnotesize
\begin{longtable}{|p{0.3\textwidth}|p{0.65\textwidth}|}
\hline
\textbf{Usuario (Rol)} & Administrador \\ \hline
\textbf{Puntos estimados} & 8 \\ \hline
\textbf{Descripción} & \textbf{Como} Administrador, \textbf{Quiero} configurar la conexión API con nuestro sistema ERP, \textbf{Para} que los datos de inventario de repuestos se sincronicen automáticamente y evitar la doble digitación.. \\ \hline
\textbf{Observaciones} & Depende de la documentación de la API del ERP externo. \\ \hline
\end{longtable}
\normalsize

\noindent
\textbf{Criterios de Aceptación:}
\footnotesize
\begin{longtable}{|p{0.25\textwidth}|p{0.2\textwidth}|p{0.2\textwidth}|p{0.25\textwidth}|}
\hline
\textbf{Escenario} & \textbf{Dado} & \textbf{Cuando} & \textbf{Entonces} \\ \hline
\endhead
Prueba de conexión API exitosa & Dado que el `Administrador' está en la pantalla de ``Integraciones'' y ha ingresado una API Key válida del ERP. & Cuando hace clic en el botón ``Probar Conexión''. & Entonces el sistema debe mostrar un mensaje de éxito: ``Conexión con ERP establecida''. \\ \hline
Prueba de conexión API fallida & Dado que el `Administrador' está en la pantalla de ``Integraciones''. & Cuando intenta probar la conexión con una API Key incorrecta o expirada. & Entonces el sistema debe mostrar un mensaje de error claro: ``Fallo la autenticación. Verifique la API Key.'' \\ \hline
\end{longtable}
\normalsize

\noindent
\textbf{Ficha Técnica: ADM-016 - Publicación de Anuncios del Sistema}
\footnotesize
\begin{longtable}{|p{0.3\textwidth}|p{0.65\textwidth}|}
\hline
\textbf{Usuario (Rol)} & Administrador \\ \hline
\textbf{Puntos estimados} & 2 \\ \hline
\textbf{Descripción} & \textbf{Como} Administrador, \textbf{Quiero} publicar un anuncio global o por rol, \textbf{Para} comunicar a los usuarios sobre ventanas de mantenimiento u otras noticias importantes de la plataforma.. \\ \hline
\textbf{Observaciones} & N/A \\ \hline
\end{longtable}
\normalsize

\noindent
\textbf{Criterios de Aceptación:}
\footnotesize
\begin{longtable}{|p{0.25\textwidth}|p{0.2\textwidth}|p{0.2\textwidth}|p{0.25\textwidth}|}
\hline
\textbf{Escenario} & \textbf{Dado} & \textbf{Cuando} & \textbf{Entonces} \\ \hline
\endhead
Publicar anuncio a un rol específico & Dado que el `Administrador' está en la pantalla de ``Publicar Anuncio''. & Cuando escribe un mensaje, selecciona el rol ``Técnico'' como destinatario y lo publica. & Entonces solo los usuarios con el rol ``Técnico'' deben ver un banner con ese anuncio la próxima vez que inicien sesión. \\ \hline
\end{longtable}
\normalsize

\noindent
\textbf{Ficha Técnica: ADM-017 - Dashboard Ejecutivo de KPIs}
\footnotesize
\begin{longtable}{|p{0.3\textwidth}|p{0.65\textwidth}|}
\hline
\textbf{Usuario (Rol)} & Gerente \\ \hline
\textbf{Puntos estimados} & 8 \\ \hline
\textbf{Descripción} & \textbf{Como} Gerente, \textbf{Quiero} visualizar un Dashboard Ejecutivo con KPIs clave (MTBF, Costos, Disponibilidad) que permita filtrar por fecha y área, \textbf{Para} identificar tendencias de fallos y ``activos problema'' (Bad Actors) para dirigir el presupuesto de mantenimiento donde más se necesita.. \\ \hline
\textbf{Observaciones} & Requiere procesamiento masivo de datos históricos y librería de gráficos. \\ \hline
\end{longtable}
\normalsize

\noindent
\textbf{Criterios de Aceptación:}
\footnotesize
\begin{longtable}{|p{0.25\textwidth}|p{0.2\textwidth}|p{0.2\textwidth}|p{0.25\textwidth}|}
\hline
\textbf{Escenario} & \textbf{Dado} & \textbf{Cuando} & \textbf{Entonces} \\ \hline
\endhead
Visualización de KPIs Globales & Dado que el `Gerente' ha iniciado sesión. & Cuando navega al módulo de ``Inteligencia de Mantenimiento''. & Entonces el sistema debe mostrar los indicadores MTBF, MTTR y Disponibilidad del \textbf{mes en curso} para toda la planta, utilizando gráficos de tendencia (líneas) y medidores (gauges). \\ \hline
Filtrado por Rango de Fecha y Zona & Dado que el `Gerente' está viendo el Dashboard General. & Cuando selecciona el rango ``Último Trimestre'' y filtra por la zona ``Casa de Bombas''. & Entonces todos los gráficos deben recalcularse automáticamente para mostrar solo los datos de esa zona en ese periodo. \\ \hline
Drill-down a Activos Críticos (Pareto) & Dado que el `Gerente' visualiza el gráfico de ``Top 10 Activos con más Paradas'' (Pareto). & Cuando hace clic en la barra del activo más crítico (el número 1). & Entonces el sistema debe redirigirlo a la ``Hoja de Vida'' de ese activo específico para ver el detalle de sus últimas fallas. \\ \hline
\end{longtable}
\normalsize

\noindent
\textbf{Ficha Técnica: ADM-018 - Exportación de Reportes en PDF}
\footnotesize
\begin{longtable}{|p{0.3\textwidth}|p{0.65\textwidth}|}
\hline
\textbf{Usuario (Rol)} & Gerente \\ \hline
\textbf{Puntos estimados} & 5 \\ \hline
\textbf{Descripción} & \textbf{Como} Gerente, \textbf{Quiero} descargar un reporte en PDF con los indicadores clave del mantenimiento (MTBF, MTTR, Disponibilidad), \textbf{Para} analizar el desempeño mensual y presentarlo en las reuniones de gestión sin depender de la conexión al sistema.. \\ \hline
\textbf{Observaciones} & Generación de documentos en servidor. Cuidar el diseño del layout PDF. \\ \hline
\end{longtable}
\normalsize

\noindent
\textbf{Criterios de Aceptación:}
\footnotesize
\begin{longtable}{|p{0.25\textwidth}|p{0.2\textwidth}|p{0.2\textwidth}|p{0.25\textwidth}|}
\hline
\textbf{Escenario} & \textbf{Dado} & \textbf{Cuando} & \textbf{Entonces} \\ \hline
\endhead
Reporte descargado correctamente & Dado que el `Gerente' está en la pantalla de “Inteligencia de Negocio”. & Cuando selecciona el mes a consultar y hace clic en “Descargar PDF”. & Entonces el sistema debe generar y descargar un archivo \textbf{.pdf} que contenga los gráficos de los indicadores seleccionados y mostrar el mensaje: “Reporte generado correctamente.” \\ \hline
Intento fallido por falta de datos & Dado que el `Gerente' intenta descargar el reporte del mes seleccionado. & Cuando el sistema detecta que no existen registros suficientes para calcular los indicadores. & Entonces debe mostrar un mensaje de error: “No hay datos disponibles para generar el reporte del mes seleccionado.” \\ \hline
\end{longtable}
\normalsize

\noindent
\textbf{Ficha Técnica: ADM-019 - Sincronización de Datos con Herramientas BI}
\footnotesize
\begin{longtable}{|p{0.3\textwidth}|p{0.65\textwidth}|}
\hline
\textbf{Usuario (Rol)} & Gerente \\ \hline
\textbf{Puntos estimados} & 5 \\ \hline
\textbf{Descripción} & \textbf{Como} Gerente, \textbf{Quiero} exportar o sincronizar los datos de los KPIs directamente con herramientas externas (como Power BI o Excel), \textbf{Para} tener presentaciones gerenciales siempre actualizadas y realizar análisis profundos sin tener que copiar y pegar gráficos manualmente cada mes.. \\ \hline
\textbf{Observaciones} & Debe exponer endpoints seguros (API REST) para consumo externo. \\ \hline
\end{longtable}
\normalsize

\noindent
\textbf{Criterios de Aceptación:}
\footnotesize
\begin{longtable}{|p{0.25\textwidth}|p{0.2\textwidth}|p{0.2\textwidth}|p{0.25\textwidth}|}
\hline
\textbf{Escenario} & \textbf{Dado} & \textbf{Cuando} & \textbf{Entonces} \\ \hline
\endhead
Exportación de Datos para BI & Dado que el `Gerente' está viendo el Dashboard de KPIs. & Cuando selecciona la opción ``Exportar Dataset (CSV/JSON)''. & Entonces el sistema debe generar un archivo estructurado y limpio, listo para ser importado en herramientas como Power BI o Excel. \\ \hline
Conexión de API de Datos & Dado que el `Gerente' utiliza una herramienta de BI externa (ej. Power BI Desktop). & Cuando configura el conector web con la URL de la API de Reportes del sistema y su credencial. & Entonces la herramienta externa debe ser capaz de importar la tabla de datos históricos de KPIs y permitir su actualización automática. \\ \hline
\end{longtable}
\normalsize

\noindent
\textbf{Ficha Técnica: ADM-022 - Sugerencias de Optimización (PMO)}
\footnotesize
\begin{longtable}{|p{0.3\textwidth}|p{0.65\textwidth}|}
\hline
\textbf{Usuario (Rol)} & Supervisor \\ \hline
\textbf{Puntos estimados} & 5 \\ \hline
\textbf{Descripción} & \textbf{Como} Supervisor, \textbf{Quiero} ver sugerencias automáticas sobre la frecuencia de mantenimiento basadas en el historial de fallas del activo, \textbf{Para} identificar equipos que están siendo sobre-mantenidos y ajustar el plan para reducir costos innecesarios. \\ \hline
\textbf{Observaciones} & Lógica de análisis de datos históricos. \\ \hline
\end{longtable}
\normalsize

\noindent
\textbf{Criterios de Aceptación:}
\footnotesize
\begin{longtable}{|p{0.25\textwidth}|p{0.2\textwidth}|p{0.2\textwidth}|p{0.25\textwidth}|}
\hline
\textbf{Escenario} & \textbf{Dado} & \textbf{Cuando} & \textbf{Entonces} \\ \hline
\endhead
\textbf{Detección de Sobre-Mantenimiento} & Dado que un activo tiene mantenimientos preventivos mensuales y cero fallas registradas en los últimos 12 meses. & Cuando el Supervisor accede al módulo de ``Optimización PMO''. & Entonces el sistema debe mostrar una tarjeta de sugerencia: ``Activo [ID] con alta confiabilidad. Se sugiere aumentar frecuencia a Bimestral''. \\ \hline
\textbf{Aplicación de Sugerencia} & Dado que el sistema muestra una sugerencia de cambio de frecuencia. & Cuando el Supervisor hace clic en ``Aplicar Cambio''. & Entonces el Plan de Mantenimiento Preventivo asociado debe actualizarse automáticamente con la nueva frecuencia y registrar quién aprobó el cambio. \\ \hline
\end{longtable}
\normalsize

\noindent
\textbf{Ficha Técnica: ADM-024 - Configuración General del Sistema}
\footnotesize
\begin{longtable}{|p{0.3\textwidth}|p{0.65\textwidth}|}
\hline
\textbf{Usuario (Rol)} & Administrador \\ \hline
\textbf{Puntos estimados} & 2 \\ \hline
\textbf{Descripción} & \textbf{Como} Administrador, \textbf{Quiero} configurar los parámetros globales del sistema (ej. Logo de la empresa, Zona Horaria, Unidades de Medida), \textbf{Para} personalizar la plataforma según la identidad y ubicación de la planta sin necesidad de modificar el código.. \\ \hline
\textbf{Observaciones} & Configuración persistente a nivel de aplicación (Singleton o DB config). \\ \hline
\end{longtable}
\normalsize

\noindent
\textbf{Criterios de Aceptación:}
\footnotesize
\begin{longtable}{|p{0.25\textwidth}|p{0.2\textwidth}|p{0.2\textwidth}|p{0.25\textwidth}|}
\hline
\textbf{Escenario} & \textbf{Dado} & \textbf{Cuando} & \textbf{Entonces} \\ \hline
\endhead
Cambio de Logo & Dado que el `Administrador' se encuentra en la pantalla de ``Configuración General''. & Cuando carga un archivo de imagen (PNG/JPG) en el campo ``Logo Corporativo'' y guarda los cambios. & Entonces el logo en la barra de navegación superior se actualiza inmediatamente para reflejar la nueva imagen. \\ \hline
Cambio de Zona Horaria & Dado que el `Administrador' está en la configuración regional y la hora actual del sistema está en UTC. & Cuando selecciona ``UTC-5 (Colombia)'' en la lista desplegable de Zona Horaria. & Entonces todas las fechas y horas mostradas en las Órdenes de Trabajo deben ajustarse automáticamente a la nueva zona seleccionada. \\ \hline
\end{longtable}
\normalsize
